% Chunk: physical PDF pages 485--491 (printed pages 462--468).
% Semantic coverage: conclusion of Section 10.7; Section 10.8 through the
% physical-page split ending with ``hadron-'' on page 491.
% Equations: (10.8.1)--(10.8.28).
% Figures: none. Footnotes: two. Uncertainties: none.

coupling constants of the composite particle. Unfortunately, it has not
been possible to implement this procedure in quantum field theories,
because as we have seen $Z=0$ means that the particle couples as strongly
as possible to its constituents, and this rules out the use of perturbation
theory. The condition $Z=0$ does prove useful in non-relativistic quantum
mechanics; for instance, it fixes the coupling of the deuteron to the neutron
and proton~\hyperref[ref:10.12]{[12]}.

Although the spectral representation has been derived here only for a
spinless field, it is easy to generalize these results to other fields. Indeed,
in the next chapter we shall show that to order $e^2$, the $Z$-factor for the
electromagnetic field (conventionally called $Z_3$) is given by
\begin{equation*}
  Z_3=1-\frac{e^2}{12\pi^2}\ln\left(\frac{\Lambda^2}{m_e^2}\right)
\end{equation*}
(where $\Lambda\gg m_e$ is an ultraviolet cutoff), in agreement with the
bound~\eqref{eq:10.7.22}.

\subsection[Dispersion Relations]{Dispersion Relations\texorpdfstring{\protect\footnotemark}{}}
\footnotetext{This section lies somewhat out of the book's main line of
development, and may be omitted in a first reading.}

The failure of early attempts to apply perturbative quantum field theory
to the strong and weak nuclear forces had led theorists by the late 1950s
to attempt the use of the analyticity and unitarity of scattering amplitudes
as a way of deriving general non-perturbative results that would not
depend on any particular field theory. This started with a revival of
interest in dispersion relations. In its original form~\hyperref[ref:10.13]{[13]},
a dispersion relation was a formula giving the real part of the index of
refraction in terms of an integral over its imaginary part. It was derived
from an analyticity property of the index of refraction as a function of
frequency, which followed from the condition that electromagnetic signals
in a medium cannot travel faster than light in a vacuum. By expressing
the index of refraction in terms of the forward photon scattering amplitude,
the dispersion relation could be rewritten as a formula for the real part of
the forward scattering amplitude as an integral of its imaginary part, and
hence via unitarity in terms of the total cross-section. One of the exciting
things about this relation was that it provided an alternative to conventional
perturbation theory; given the scattering amplitude to order $e^2$, one could
calculate the cross-section and the imaginary part of the scattering amplitude
to order $e^4$, and then use the dispersion relation to calculate the real
part of the forward scattering amplitude to this order, without having ever
to calculate a loop graph.

The modern approach to dispersion relations began in 1954 with the work
of Gell-Mann, Goldberger, and Thirring~\hyperref[ref:10.14]{[14]}. Instead
of considering the propagation of light in a medium, they derived the
analyticity of the scattering amplitude directly from the condition of
microscopic causality, which states that commutators of field operators vanish
when the points at which the operators are evaluated are separated by a
space-like interval. This approach allowed Goldberger~\hyperref[ref:10.15]{[15]}
soon thereafter to derive a very useful dispersion relation for the forward
pion--nucleon scattering amplitude.

To see how to use the principle of microscopic causality, consider the
forward scattering in the laboratory frame of a massless boson of any spin on
an arbitrary target $\alpha$ of mass $m_\alpha>0$ and $\mathbf p_\alpha=0$.
(This has important applications to the scattering not only of photons but
also pions in the limit $m_\pi=0$, to be discussed in Volume II.) By a repeated
use of Eq.~\eqref{eq:10.3.4} or the Lehmann--Symanzik--Zimmermann
theorem~\hyperref[ref:10.3]{[3]}, the $S$-matrix element here is
\begin{align}
  S={}&\frac{1}{(2\pi)^3\sqrt{4\omega\omega'}|N|^2}
  \lim_{k^2\to0}\lim_{k'^2\to0}\notag\\
  &\quad\cdot\int d^4x\int d^4y\,e^{-ik'\cdot y}e^{ik\cdot x}
  (i\Box_y)(i\Box_x)
  \bra{\alpha}T\{A^\dagger(y),A(x)\}\ket{\alpha}.
  \tag{10.8.1}\label{eq:10.8.1}
\end{align}
Here $k$ and $k'$ are the initial and final boson four-momenta, with
$\omega=k^0$, $\omega'=k'^0$; $A(x)$ is any Heisenberg-picture operator
with a non-vanishing matrix element
$\bra{\Omega}A(x)\ket{k}=(2\pi)^{-3/2}(2\omega)^{-1/2}Ne^{ik\cdot x}$
between the one-boson state $\ket{k}$ and the vacuum; and $N$ is the
constant in this matrix element. In photon scattering $A(x)$ would be one
of the transverse components of the electromagnetic field, while for massless
pion scattering it would be a pseudoscalar function of hadron fields. The
differential operators $-i\Box_x$ and $-i\Box_y$ are inserted to supply
factors of $ik'^2$ and $ik^2$ that are needed to cancel the external line
boson propagators. Letting these operators act on $A^\dagger(y)$ and $A(x)$,
we have
\begin{align}
  S={}&\frac{-1}{(2\pi)^3\sqrt{4\omega\omega'}|N|^2}
  \lim_{k^2\to0}\lim_{k'^2\to0}\notag\\
  &\quad\cdot\int d^4x\int d^4y\,e^{-ik'\cdot y}e^{ik\cdot x}
  \bra{\alpha}T\{J^\dagger(y),J(x)\}\ket{\alpha}+\mathrm{ETC},
  \tag{10.8.2}\label{eq:10.8.2}
\end{align}
where $J(x)\equiv\Box_xA(x)$, and `ETC' denotes the Fourier transform
of equal time commutator terms arising from the derivative acting on the
step functions in the time-ordered product. The commutators of operators
like $A(x)$ and $A^\dagger(y)$ (or their derivatives) vanish for $x^0=y^0$
unless $\mathbf{x}=\mathbf{y}$, so the `ETC' term is the Fourier transform
of a differential operator acting on $\delta^4(x-y)$, and is hence a
polynomial function of the boson four-momenta. We are concerned here with
the analytic properties of the $S$-matrix element, so the details of this
polynomial will be irrelevant.

Using translation invariance, Eq.~\eqref{eq:10.8.2} gives the $S$-matrix
element as $S=-2\pi i\delta^4(k'-k)M(\omega)$, where
\begin{equation}
  M(\omega)=\frac{-i}{2\omega|N|^2}F(\omega),
  \tag{10.8.3}\label{eq:10.8.3}
\end{equation}
\begin{equation}
  F(\omega)\equiv\int d^4x\,e^{i\omega\ell\cdot x}
  \bra{\alpha}T\{J^\dagger(0),J(x)\}\ket{\alpha}+\mathrm{ETC},
  \tag{10.8.4}\label{eq:10.8.4}
\end{equation}
it now being understood that $k^\mu=\omega\ell^\mu$, where $\ell$ is a
fixed four-vector with $\ell^\mu\ell_\mu=0$ and $\ell^0=1$.

The time-ordered product can be rewritten in terms of commutators in two
different ways:
\begin{align}
  T\{J^\dagger(0),J(x)\}
  &=\theta(-x^0)[J^\dagger(0),J(x)]+J(x)J^\dagger(0)\notag\\
  &=-\theta(x^0)[J^\dagger(0),J(x)]+J^\dagger(0)J(x).
  \tag{10.8.5}\label{eq:10.8.5}
\end{align}
Correspondingly, we can write
\begin{equation}
  F(\omega)=F_A(\omega)+F_+(\omega)=F_R(\omega)+F_-(\omega),
  \tag{10.8.6}\label{eq:10.8.6}
\end{equation}
where
\begin{equation}
  F_A(\omega)\equiv\int d^4x\,\theta(-x^0)
  \bra{\alpha}[J^\dagger(0),J(x)]\ket{\alpha}
  e^{i\omega\ell\cdot x}+\mathrm{ETC},
  \tag{10.8.7}\label{eq:10.8.7}
\end{equation}
\begin{equation}
  F_R(\omega)\equiv-\int d^4x\,\theta(x^0)
  \bra{\alpha}[J^\dagger(0),J(x)]\ket{\alpha}
  e^{i\omega\ell\cdot x}+\mathrm{ETC},
  \tag{10.8.8}\label{eq:10.8.8}
\end{equation}
\begin{equation}
  F_+(\omega)\equiv\int d^4x\,
  \bra{\alpha}J(x)J^\dagger(0)\ket{\alpha}e^{i\omega\ell\cdot x},
  \tag{10.8.9}\label{eq:10.8.9}
\end{equation}
\begin{equation}
  F_-(\omega)\equiv\int d^4x\,
  \bra{\alpha}J^\dagger(0)J(x)\ket{\alpha}e^{i\omega\ell\cdot x}.
  \tag{10.8.10}\label{eq:10.8.10}
\end{equation}
Microscopic causality tells us that the integrands in~\eqref{eq:10.8.7}
and~\eqref{eq:10.8.8} vanish unless $x^\mu$ is within the light cone, and
the step functions then require that $x^\mu$ is in the backward light cone
in~\eqref{eq:10.8.7}, so that $x\cdot\ell>0$, and in the forward light cone
in Eq.~\eqref{eq:10.8.8}, so that $x\cdot\ell<0$. We conclude that
$F_A(\omega)$ is analytic for $\operatorname{Im}\omega>0$ and $F_R(\omega)$
is analytic for $\operatorname{Im}\omega<0$, because in both cases the factor
$e^{i\omega\ell\cdot x}$ provides a cutoff for the integral over $x^\mu$.
(Recall that the `ETC' term is a polynomial, and hence analytic at all finite
points.) We may then define a function
\begin{equation}
  \mathcal F(\omega)\equiv
  \begin{cases}
    F_A(\omega),&\operatorname{Im}\omega>0,\\
    F_R(\omega),&\operatorname{Im}\omega<0,
  \end{cases}
  \tag{10.8.11}\label{eq:10.8.11}
\end{equation}
which is analytic in the whole complex $\omega$ plane, except for a cut on
the real axis.

We can now derive the dispersion relation. According to
Eq.~\eqref{eq:10.8.6}, the discontinuity of $\mathcal F(\omega)$ across the
cut at any real $E$ is
\begin{equation}
  \mathcal F(E+i\epsilon)-\mathcal F(E-i\epsilon)
  =F_A(E)-F_R(E)=F_-(E)-F_+(E).
  \tag{10.8.12}\label{eq:10.8.12}
\end{equation}
If $\mathcal F(\omega)/\omega^n$ vanishes as $|\omega|\to\infty$ in the
upper or lower half-plane, then by dividing by any polynomial $P(\omega)$
of order $n$ we obtain a function that vanishes for $|\omega|\to\infty$
and is analytic except for the cut on the real axis and poles at the zeroes
$\omega_\nu$ of $P(\omega)$. (Where $\mathcal F(\omega)$ itself vanishes
as $|\omega|\to\infty$, we can take $P(\omega)=1$.) According to the
method of residues, we then have
\begin{equation}
  \frac{\mathcal F(\omega)}{P(\omega)}
  +\sum_\nu\frac{\mathcal F(\omega_\nu)}
  {(\omega_\nu-\omega)P'(\omega_\nu)}
  =\frac{1}{2\pi i}\oint_C
  \frac{\mathcal F(z)\,dz}{(z-\omega)P(z)},
  \tag{10.8.13}\label{eq:10.8.13}
\end{equation}
where $\omega$ is any point off the real axis, and $C$ is a contour consisting
of two segments: one running just above the real axis from $-\infty+i\epsilon$
to $+\infty+i\epsilon$ and then around a large semi-circle back to
$-\infty+i\epsilon$, and the other just below the real axis from
$+\infty-i\epsilon$ to $-\infty-i\epsilon$ and then around a large
semi-circle back to $+\infty-i\epsilon$. Because the function
$\mathcal F(z)/P(z)$ vanishes for $|z|\to\infty$, we can neglect the
contribution from the large semi-circles. Using Eq.~\eqref{eq:10.8.12},
Eq.~\eqref{eq:10.8.13} becomes
\begin{equation}
  \mathcal F(\omega)=Q(\omega)+\frac{P(\omega)}{2\pi i}
  \int_{-\infty}^{+\infty}
  \frac{F_-(E)-F_+(E)}{(E-\omega)P(E)}\,dE,
  \tag{10.8.14}\label{eq:10.8.14}
\end{equation}
where $Q(\omega)$ is the $(n-1)$th-order polynomial
\begin{equation*}
  Q(\omega)=-P(\omega)\sum_\nu
  \frac{\mathcal F(\omega_\nu)}
  {(\omega_\nu-\omega)P'(\omega_\nu)}.
\end{equation*}
A dispersion relation of this form, with $P(\omega)$ and $Q(\omega)$ of
order $n$ and $n-1$ respectively, is said to have $n$ \emph{subtractions}.
If we can take $P=1$ then $Q=0$, and the dispersion relation is said to be
unsubtracted.

If we now let $\omega$ approach the real axis from above,
Eq.~\eqref{eq:10.8.14} gives
\begin{equation}
  F_A(\omega)=Q(\omega)+\frac{P(\omega)}{2\pi i}
  \int_{-\infty}^{+\infty}
  \frac{F_-(E)-F_+(E)}{(E-\omega-i\epsilon)P(E)}\,dE.
  \tag{10.8.15}\label{eq:10.8.15}
\end{equation}
Recalling Eqs.~\eqref{eq:10.8.6} and~\eqref{eq:3.1.25}, this is
\begin{align}
  F(\omega)={}&Q(\omega)+\frac12F_-(\omega)+\frac12F_+(\omega)
  \notag\\
  &{}+\frac{P(\omega)}{2\pi i}\int_{-\infty}^{+\infty}
  \frac{F_-(E)-F_+(E)}{(E-\omega)P(E)}\,dE
  \tag{10.8.16}\label{eq:10.8.16}
\end{align}
with $1/(E-\omega)$ now interpreted as the principal value function
$\mathcal P/(E-\omega)$.

This result is useful because the functions $F_\pm(E)$ may be expressed
in terms of measurable cross-sections. Summing over a complete set of
multi-particle intermediate states $\beta$ in Eqs.~\eqref{eq:10.8.9}
and~\eqref{eq:10.8.10} (including integrations over the momenta of the
particles in $\beta$) and using translation invariance again, we have
\begin{equation}
  F_+(E)=(2\pi)^4\sum_\beta
  \left|\bra{\beta}J(0)^\dagger\ket{\alpha}\right|^2
  \delta^4(-p_\alpha+E\ell+p_\beta),
  \tag{10.8.17}\label{eq:10.8.17}
\end{equation}
\begin{equation}
  F_-(E)=(2\pi)^4\sum_\beta
  \left|\bra{\beta}J(0)\ket{\alpha}\right|^2
  \delta^4(p_\alpha+E\ell-p_\beta).
  \tag{10.8.18}\label{eq:10.8.18}
\end{equation}
But the matrix elements for the absorption of the massless scalar boson
$B$ in $B+\alpha\to\beta$ or its antiparticle $B^c$ in
$B^c+\alpha\to\beta$ are
\begin{equation}
  -2i\pi M_{B^c+\alpha\to\beta}
  =\frac{(2\pi)^4}{(2\pi)^{3/2}\sqrt{2E_{B^c}}N}
  \bra{\beta}J^\dagger(0)\ket{\alpha},
  \tag{10.8.19}\label{eq:10.8.19}
\end{equation}
\begin{equation}
  -2i\pi M_{B+\alpha\to\beta}
  =\frac{(2\pi)^4}{(2\pi)^{3/2}\sqrt{2E_B}N}
  \bra{\beta}J(0)\ket{\alpha}.
  \tag{10.8.20}\label{eq:10.8.20}
\end{equation}
Comparing with Eq.~\eqref{eq:3.4.15}, we see that $F_\pm(E)$ may be
expressed in terms of total cross-sections\footnote{In some cases where
selection rules allow the transition $\alpha\to\alpha+B$ and
$\alpha\to\alpha+B^c$, the functions $F_\pm(E)$ also contain terms
proportional to $\delta(E)$ arising from the contribution of the one-particle
state $\alpha$ in the sum over intermediate states $\beta$. This does not
occur for transversely polarized photons, or for pseudoscalar pions in the
limit $m_\pi\to0$.} at energies $\mp E$:
\begin{equation}
  F_+(E)=\theta(-E)\frac{2|E||N|^2}{(2\pi)^3}
  \sigma_{\alpha+B^c}(|E|),
  \tag{10.8.21}\label{eq:10.8.21}
\end{equation}
\begin{equation}
  F_-(E)=\theta(E)\frac{2E|N|^2}{(2\pi)^3}
  \sigma_{\alpha+B}(E).
  \tag{10.8.22}\label{eq:10.8.22}
\end{equation}
The scattering amplitude~\eqref{eq:10.8.3} is now, for real $\omega>0$,
\begin{align}
  M(\omega)={}&\frac{-iQ(\omega)}{2\omega|N|^2}
  -\frac{i}{2(2\pi)^3}\sigma_{\alpha+B}(\omega)\notag\\
  &{}-\frac{P(\omega)}{\omega(2\pi)^4}\int_0^\infty
  \left[\frac{\sigma_{\alpha+B}(E)}{(E-\omega)P(E)}
  +\frac{\sigma_{\alpha+B^c}(E)}{(E+\omega)P(-E)}\right]E\,dE.
  \tag{10.8.23}\label{eq:10.8.23}
\end{align}
It is more usual to express this dispersion relation in terms of the
amplitude $f(\omega)$ for forward scattering in the laboratory frame,
defined so that the laboratory frame differential cross-section in the forward
direction is $|f(\omega)|^2$. This amplitude is given in terms of $M(\omega)$
by $f(\omega)=-4\pi^2\omega M(\omega)=2\pi^2iF(\omega)/|N|^2$, so
Eq.~\eqref{eq:10.8.23} now reads
\begin{align*}
  f(\omega)={}&R(\omega)+\frac{i\omega}{4\pi}
  \sigma_{\alpha+B}(\omega)\\
  &{}+\frac{P(\omega)}{4\pi^2}\int_0^\infty
  \left[\frac{\sigma_{\alpha+B}(E)}{(E-\omega)P(E)}
  +\frac{\sigma_{\alpha+B^c}(E)}{(E+\omega)P(-E)}\right]E\,dE,
\end{align*}
where $R(\omega)\equiv2i\pi^2Q(\omega)/|N|^2$. The optical theorem
\eqref{eq:3.6.4} tells us that the second term on the right-hand side equals
$i\operatorname{Im}f(\omega)$, so this can just as well be written in the more
conventional form
\begin{align}
  \operatorname{Re}f(\omega)={}&R(\omega)
  +\frac{P(\omega)}{4\pi^2}\int_0^\infty
  \left[\frac{\sigma_{\alpha+B}(E)}{(E-\omega)P(E)}
  +\frac{\sigma_{\alpha+B^c}(E)}{(E+\omega)P(-E)}\right]E\,dE.
  \tag{10.8.24}\label{eq:10.8.24}
\end{align}
In particular, we see that $R(\omega)$ is real if we choose $P(\omega)$ real.

The forward scattering amplitude also satisfies an important symmetry
condition. By changing the integration variable in Eqs.~\eqref{eq:10.8.7}
and~\eqref{eq:10.8.8} from $x$ to $-x$ and then using the
translation-invariance property
\begin{equation*}
  \bra{\alpha}[J^\dagger(0),J(-x)]\ket{\alpha}
  =\bra{\alpha}[J^\dagger(x),J(0)]\ket{\alpha}
\end{equation*}
we see that for $\operatorname{Im}\omega\leq0$, $F_A(-\omega)$ is the same
as $F_R(\omega)$, except for an interchange of $J$ with $J^\dagger$. That is,
\begin{equation*}
  F_A(-\omega)=F_R^c(\omega)
  \qquad\text{for}\qquad\operatorname{Im}\omega\leq0,
\end{equation*}
where a superscript $c$ indicates that the amplitude refers to the scattering
of the antiparticle $B^c$ on $\alpha$. (We leave it to the reader to show
that this relation is not upset by the equal-time commutator terms in
Eqs.~\eqref{eq:10.8.7} and~\eqref{eq:10.8.8}.) In the same way, we find
\begin{equation*}
  F_R(-\omega)=F_A^c(\omega)
  \qquad\text{for}\qquad\operatorname{Im}\omega\geq0,
\end{equation*}
and for real $\omega$
\begin{equation*}
  F_\pm(-\omega)=F_\mp^c(\omega).
\end{equation*}
Using these relations in~\eqref{eq:10.8.6}, and recalling that $f(\omega)$ is
proportional to $F(\omega)$, we find the \emph{crossing symmetry} relation,
that for real $\omega$
\begin{equation}
  f(-\omega)=f^c(\omega).
  \tag{10.8.25}\label{eq:10.8.25}
\end{equation}

We are free to take $P(\omega)$ as any polynomial of sufficiently high order,
but $R(\omega)$ then depends not only on $P(\omega)$ but also on the values
of $\mathcal F(\omega)$ at the zeroes of $P(\omega)$. For $P(\omega)$ real
and of $n$th order, the only free parameters in Eq.~\eqref{eq:10.8.16} are
the $n$ real coefficients in the real $(n-1)$th order polynomial $R(\omega)$.
Hence Eq.~\eqref{eq:10.8.16} contains just $n$ unknown real independent
constants, the coefficients in the polynomial $R(\omega)$ for a given
$P(\omega)$. We therefore wish to take the order $n$ of the otherwise
arbitrary polynomial $P(\omega)$ to be as small as possible.

We might try taking $P(\omega)=1$, but this doesn't work. The analysis of
Section 3.7 suggests that the forward scattering amplitude should grow like
$\omega$ or perhaps as fast as $\omega\ln^2\omega$. In this case for
$f(\omega)/P(\omega)$ to vanish as $\omega\to0$, it is sufficient to take
$P(\omega)$ as a second order polynomial, so that $R(\omega)$ is linear in
$\omega$. Choosing $P(E)=E^2$ for convenience, Eq.~\eqref{eq:10.8.24}
then becomes
\begin{align}
  \operatorname{Re}f(\omega)={}&a+b\omega
  +\frac{\omega^2}{4\pi^2}\int_0^\infty
  \left[\frac{\sigma_{\alpha+B}(E)}{E-\omega}
  +\frac{\sigma_{\alpha+B^c}(E)}{E+\omega}\right]\frac{dE}{E},
  \tag{10.8.26}\label{eq:10.8.26}
\end{align}
with $a$ and $b$ unknown real constants. The crossing symmetry condition
\eqref{eq:10.8.25} tells us that the corresponding constants in the dispersion
relation for the antiparticle scattering amplitude $f^c(\omega)$ are
\begin{equation}
  a^c=a,\qquad b^c=-b.
  \tag{10.8.27}\label{eq:10.8.27}
\end{equation}
If we assume for instance that the cross-sections $\sigma_{\alpha+B}(E)$ and
$\sigma_{\alpha+B^c}(E)$ behave for $E\to\infty$ as different constants
times $(\ln E)^r$, then~\eqref{eq:10.8.26} would give
\begin{equation}
  \operatorname{Re}f(\omega)\sim
  [\sigma_{\alpha+B}(\omega)-\sigma_{\alpha+B^c}(\omega)]\ln\omega
  \sim\omega(\ln\omega)^{r+1}
  \tag{10.8.28}\label{eq:10.8.28}
\end{equation}
so the real part of the scattering amplitude would grow faster than the
imaginary part by a factor $\ln\omega$. This is implausible; we saw in Section
3.7 that the real part of the forward scattering amplitude is expected to
become much smaller than the imaginary part for $\omega\to\infty$, as
confirmed by experiment. We conclude that if $\sigma_{\alpha+B}(E)$ and
$\sigma_{\alpha+B^c}(E)$ do behave for $E\to\infty$ as constants times
$(\ln E)^r$ then the constants must be the same. Because we are concerned
here with the high-energy limit, this result does not depend on the assumption
that $B$ is a massless boson, so in the same sense, \emph{the ratio of the
cross-sections of any particle and its antiparticle on a fixed target should
approach unity at high energy}. This result is a somewhat generalized version
of what is known as Pomeranchuk's theorem~\hyperref[ref:10.16]{[16]}.
(Pomeranchuk considered only the case $r=0$, while Section 3.7 and the
observed behavior of cross-sections both suggest that $r=2$ is more likely.)

Although Pomeranchuk took his estimates of the asymptotic behavior of
scattering amplitudes from arguments like those of Section 3.7, today high
energy behavior is usually inferred from Regge pole theory~\hyperref[ref:10.17]{[17]}.
It would take us too far from our subject to go into details about this;
suffice it to say that for hadronic processes the asymptotic behavior of
$f(\omega)$ as $\omega$ goes to infinity is a sum over terms proportional to
$\omega^{\alpha_n(0)}$, where $\alpha_n(t)$ are a set of `Regge trajectories',
each representing the exchange of an infinite family of different one-hadron
states in the collision process. The leading trajectory (actually, a complex
of many trajectories) in
